\chapter{KQL Databases and Eventhouses}
\index{Eventhouse}\index{KQL}\index{Kusto Query Language}\index{materialized views}\index{arg\_max}\index{time series}\index{anomaly detection}\index{make-series}\index{late-arriving data}\index{stream joins}%

\begin{objectives}
\begin{itemize}
    \item Explain the Eventhouse architecture and how KQL databases store and index streaming data.
    \item Write core KQL: \texttt{search}, \texttt{where}, \texttt{project}, \texttt{summarize}, and time-bounded filters.
    \item Maintain incremental state with materialized views and deduplicate replayed events with \texttt{arg\_max}.
    \item Build time series with \texttt{make-series}, rolling windows, and \texttt{series\_decompose\_anomalies}.
    \item Express stream-to-stream and stream-to-dimension joins with time-windowed keys, and reason about late arrivals.
    \item Troubleshoot a dashboard whose daily aggregates are wrong because offline devices uploaded data late.
\end{itemize}
\end{objectives}

\begin{crosscloud}
Purpose-built time-series/log analytics engines---columnar, append-optimized, query-language-native---exist across the ecosystem. Fabric's Eventhouse is the Azure Data Explorer (Kusto) lineage delivered inside the SaaS workspace.

\vspace{2mm}
{\footnotesize
\begin{tabular}{@{}p{2.8cm}p{3.2cm}p{3.2cm}p{3.2cm}p{3.2cm}@{}}
\toprule
\textbf{Capability} & \textbf{Fabric} & \textbf{AWS} & \textbf{GCP} & \textbf{Snowflake} \\
\midrule
Telemetry/log engine & Eventhouse (KQL database) & CloudWatch Logs Insights / OpenSearch / Timestream & Cloud Logging / BigQuery streaming analytics & Snowflake tables + time-series functions \\
Query language & KQL & Logs Insights QL / SQL & SQL / Logging query language & SQL \\
Time-series functions & \texttt{make-series}, \texttt{series\_*} family & Timestream SQL functions & BigQuery window/ML functions & Window functions, ASOF joins \\
Incremental state & Materialized views & Scheduled queries / rollups & BigQuery materialized views & Materialized views, dynamic tables \\
Dedup idiom & \texttt{arg\_max} by id & Window dedup & \texttt{ROW\_NUMBER} dedup & \texttt{ROW\_NUMBER} / \texttt{MERGE} \\
Retention control & Retention policies per table & Log group retention / TTL & Log bucket retention / table TTL & Time Travel + data retention \\
\bottomrule
\end{tabular}}

\vspace{1mm}
Databricks approaches telemetry with Structured Streaming plus Delta; MongoDB with time-series collections. The portable skills: think in \emph{append-only streams plus incremental state}, bound your lookback windows, and treat late arrivals as a policy decision, not a surprise.
\end{crosscloud}

\begin{keyterms}
\textbf{Eventhouse}: the Fabric item that hosts one or more KQL databases for real-time analytics. \quad
\textbf{KQL (Kusto Query Language)}: a pipe-composed, read-optimized query language for telemetry and time-series data. \quad
\textbf{Materialized view}: a continuously maintained query result---incremental state over an append-only stream. \quad
\textbf{\texttt{arg\_max}}: the KQL aggregate that returns the row with the maximum value of an expression---the dedup idiom. \quad
\textbf{\texttt{make-series}}: the operator that converts event streams into ordered time-series arrays for analysis. \quad
\textbf{Late arrival}: an event whose event time precedes its arrival by more than the processing window allows.
\end{keyterms}

\section{Week 14 Roadmap}
Week~14 teaches the query heart of real-time intelligence. Monday covers Eventhouses and KQL databases. Tuesday builds KQL fluency---\texttt{search}, \texttt{where}, \texttt{project}, \texttt{summarize}---plus materialized views and the \texttt{arg\_max} dedup idiom. Wednesday goes advanced: time-series analysis, stream joins with time-windowed keys, late-arrival implications, and why unbounded lookback windows degrade performance. Thursday is the hands-on project: a five-minute rolling average and anomaly detection over the Week~13 IoT stream. Friday runs the troubleshooting scenario---a dashboard wrong because devices went offline and uploaded late---and closes the loop on end-to-end exactly-once reasoning.

KQL is the highest-leverage new language in this program. It reads like a pipeline---data flows through \texttt{|} stages---and it is optimized for exactly the questions streaming systems ask: filter by time, aggregate by window, detect the anomaly, find the latest state \cite{microsoftKQL,microsoftFabricEventhouse}.

\begin{notebox}
The unifying idea of the week: an append-only event stream plus maintained incremental state (materialized views) reproduces most of what mutable tables give you---without giving up the replay-safety of never updating in place.
\end{notebox}

\section{Monday: Eventhouses and KQL Databases}
\index{Eventhouse!architecture}\index{KQL database}
An Eventhouse is the Fabric container for KQL databases: the compute and storage fabric that ingests streams, indexes columns, compresses extents, and serves queries within seconds of ingestion \cite{microsoftFabricEventhouse}. A KQL database inside it holds tables, materialized views, functions, and policies---retention, caching, row-level security.

\begin{definitionbox}
An \textbf{Eventhouse} is Fabric's real-time analytics engine host: it provides the compute and storage under which KQL databases ingest, index, and serve high-volume append-only telemetry with sub-second-to-second query latency.
\end{definitionbox}

Three properties shape engineering practice. \emph{Append-first}: tables are optimized for relentless inserts; updates and deletes exist but are exceptions, which is why dedup happens at query/materialization time. \emph{Time-native}: every table has an ingestion timestamp, most have an event timestamp, and the engine indexes time aggressively. \emph{Policy-driven}: retention and caching are table policies, so hot recent data serves from cache while history ages out deliberately.

\begin{pitfall}
Eventhouse compute meters against your Fabric capacity. An unbounded dashboard query (``all events, all time, every refresh'') is a capacity incident disguised as a visualization. Bound every query with a time filter as a habit, not an optimization.
\end{pitfall}

\section{Tuesday: KQL Basics, Materialized Views, and Dedup}
\index{KQL!basics}\index{materialized views}\index{arg\_max!dedup}
\subsection{The Core Verbs}
KQL composes operators left to right through the pipe \cite{microsoftKQL}:

\begin{lstlisting}[language=SQL,caption={Core KQL verbs over the IoT stream.},label={lst:chapter14-core}]
-- search: raw text/expression hunt across columns
IoTTelemetry
| search "device-07";

-- where + project: filter and shape
IoTTelemetry
| where Timestamp > ago(1h) and Temperature > 25
| project DeviceId, Timestamp, Temperature;

-- summarize: aggregate by groups and time bins
IoTTelemetry
| where Timestamp > ago(24h)
| summarize AvgTemp = avg(Temperature), Readings = count()
    by DeviceId, bin(Timestamp, 5m)
| order by DeviceId, bin_Timestamp asc;
\end{lstlisting}

\subsection{Materialized Views as Incremental State}
A materialized view maintains an aggregate or projection continuously as data arrives---the streaming answer to ``recompute the summary every night.'' Over an append-only stream, the view is the state: dashboards and rules read the view, never the raw firehose \cite{microsoftKQLMV}.

\subsection{Dedup with \texttt{arg\_max}}
Chapter~13's at-least-once reality lands here. The dedup idiom keeps the latest arrival per logical event:

\begin{lstlisting}[language=SQL,caption={Deduplication idiom: latest arrival per event id.},label={lst:chapter14-dedup}]
IoTTelemetry
| summarize arg_max(ingestion_time(), *) by EventId
\end{lstlisting}

Enforced inside a materialized view (or as the first stage of every consumer query), this makes replay harmless: re-delivered events collapse back to one logical row. The dedup key---\texttt{EventId} stamped at the producer---is the contract from Chapter~13 made executable.

\section{Wednesday: Time Series, Joins, and Late Arrivals}
\index{time series}\index{make-series}\index{series\_decompose\_anomalies}\index{stream joins}\index{late-arriving data}
\subsection{Time-Series Analysis}
\texttt{make-series} converts events into ordered arrays per entity; the \texttt{series\_*} function family then operates on the arrays---filters, decomposition, forecasting, anomaly detection \cite{microsoftKQL}.

\subsection{Stream Joins with Time-Windowed Keys}
Stream-to-stream joins correlate two unbounded flows; the join key alone is insufficient---a time window bounds which events can match. Stream-to-dimension joins enrich events with the current (or as-of) dimension state. Both patterns appear in Friday's listing.

\subsection{Late Arrivals and Window Policy}
When a device goes offline and backhauls hours of events, every windowed aggregate over the affected period was computed on incomplete data. The engineering responses, in order of preference: \emph{size windows and materialization lookback to the realistic lateness distribution}; \emph{re-aggregate on arrival} where the materialized view semantics allow; and \emph{document the lateness policy} so ``wrong'' numbers are understood as a stated contract, not a defect. Unbounded lookback---``just re-scan everything always''---degrades both query and materialization performance precisely because it abandons the incremental-state design.

\begin{pitfall}
A dashboard that silently recomputes with late data shows one number at 09:00 and a different number at 17:00 for the same day, and someone will present both in the same meeting. Decide and publish the lateness policy: how late counts, when numbers stabilize, and how revisions are marked.
\end{pitfall}

\section{Thursday: Hands-on Project}
\index{hands-on lab}\index{anomaly detection!hands-on}
The Week~14 lab queries the Week~13 IoT data: a five-minute rolling average temperature per device and anomaly detection using built-in series functions.

\subsection{Goal and Architecture}
Reuse \texttt{kqldb\_iot.IoTTelemetry} from Week~13 (or regenerate with its producer). Deliverables: the rolling-average query, the anomaly-flagged result set, and a materialized view maintaining the per-device five-minute means.

\subsection{Step-by-Step Actions}
\begin{enumerate}
    \item Verify data is flowing (or restart the producer) and bound every query with \texttt{ago()}.
    \item Build the \texttt{make-series} base: average temperature per device per five-minute step.
    \item Apply \texttt{series\_fir} for the rolling average.
    \item Apply \texttt{series\_decompose\_anomalies} with a 1.5 threshold and expand back to rows.
    \item Validate: inject a synthetic hot streak from the producer and watch it flag.
    \item Create the materialized view for the per-device means and compare its latency against the raw query.
\end{enumerate}

\begin{lstlisting}[language=SQL,caption={Week 14 lab: rolling average and anomaly detection per device.},label={lst:chapter14-lab}]
// 5-minute rolling average + anomaly detection per device
IoTTelemetry
| where Timestamp > ago(6h)
| make-series AvgTemp = avg(Temperature) default=0
              on Timestamp step 5m by DeviceId
| extend RollingAvg = series_fir(AvgTemp, repeat(1, 5))
| extend Anomalies = series_decompose_anomalies(RollingAvg, 1.5)
| mv-expand Timestamp, RollingAvg, Anomalies
| project DeviceId,
          Timestamp = todatetime(Timestamp),
          RollingAvg = todouble(RollingAvg),
          IsAnomaly = toint(Anomalies) == 1
| order by DeviceId, Timestamp
\end{lstlisting}

\subsection{Validation Checklist}
\begin{itemize}
    \item Every query carries a bounded time filter.
    \item The rolling average smooths a noisy device without hiding the injected hot streak.
    \item Anomaly flags fire on the synthetic spike and stay quiet on normal noise.
    \item The materialized view returns per-device means faster than the raw aggregation.
    \item You can explain what \texttt{default=0} does to sparse series---and when that default lies.
\end{itemize}

\section{Friday: Use Case and Architecture}
\index{late-arriving data!troubleshooting}\index{exactly-once}
A streaming dashboard shows daily aggregations that keep changing: mobile devices went offline overnight and uploaded their buffered events at dawn. Daily totals for ``yesterday'' grew between the morning standup and the afternoon review. Diagnose and fix.

\subsection{Diagnosis}
The aggregates are computed over arrival-time windows, not event-time windows with a lateness policy. Yesterday's numbers were computed from whatever had \emph{arrived} by midnight, then silently revised as the backhaul landed. Three compounding issues: no event-time windowing, no stated lateness policy, and---discovered during the fix---duplicates from at-least-once redelivery inflating the revised totals.

\subsection{The Fix: End-to-End Exactly-Once as a Discipline}
Assemble the week's concepts into the full pattern: Eventstream's at-least-once delivery; idempotent ingestion into KQL and the lakehouse; the dedup key stamped at the producer and enforced at the first queryable boundary; and consumer checkpointing so restarts resume rather than replay. The query layer implements it:

\begin{lstlisting}[language=SQL,caption={Dedup, windowed stream join, and dimension enrichment in one pattern.},label={lst:chapter14-friday}]
// Dedup (keep latest arrival per id), windowed stream join, dimension lookup
let Telemetry =
    IoTTelemetry
    | where Timestamp > ago(1h)              // bounded lookback window
    | summarize arg_max(ingestion_time(), *) by EventId;
let Alerts =
    DeviceAlerts
    | where AlertTime > ago(1h)
    | summarize arg_max(ingestion_time(), *) by AlertId;
Telemetry
| join kind=inner Alerts on DeviceId
| where abs(datetime_diff('second', Timestamp, AlertTime)) <= 300
| join kind=inner (                          // stream-to-dimension
    Devices | project DeviceId, Site, Model
) on DeviceId
| summarize Readings = count() by Site, Model, bin(Timestamp, 1m)
\end{lstlisting}

For the dashboard itself: aggregate on event time, set the materialization lookback to the measured lateness distribution (for example 24 hours for this fleet), mark values inside the lateness window as \emph{provisional} in the report, and let the materialized view absorb revisions instead of the executives.

\begin{pitfall}
Eventstream replay is the top source of silent duplicates downstream: after a consumer failure or reconfiguration, the source's retention window redelivers, and every at-least-once hop can double-land rows. Stamp the dedup key at the producer, propagate it untouched, and enforce dedup at the first queryable boundary so every downstream consumer---dashboards, Reflex rules, reverse ETL---reads already-clean data.
\end{pitfall}

\section{Operational Guidance for Week 14}
KQL estates stay healthy through policy: retention, caching, bounded queries, and materialized views doing the heavy lifting dashboards would otherwise do badly.

\subsection{Performance and Reliability}
Bound every query with time. Push repeated aggregations into materialized views. Keep lookback windows honest about lateness. Use \texttt{ingestion\_time()} deliberately: it measures the pipeline, not the world. Review extent sizes and caching policies monthly against query patterns.

\subsection{Security and Governance Checklist}
\begin{itemize}
    \item Apply row-level security policies on KQL tables where telemetry carries tenant or personal scoping.
    \item Retention policies reflect both cost and compliance---telemetry is not free to keep forever.
    \item Restrict who can alter materialized views; they are production state.
    \item Classify telemetry tables; device identifiers joined to customers are personal data.
    \item Record the lateness policy per dashboard as a published contract.
    \item Audit \texttt{arg\_max} dedup usage: consumers bypassing the dedup view read dirty data.
\end{itemize}

\section{Interview Q\&A}
\begin{enumerate}
    \item \textbf{Q: How does a KQL materialized view maintain incremental state over an append-only stream?}\\
    A: It continuously applies its aggregation or projection to newly arrived extents, so consumers read maintained state instead of re-scanning history---the streaming replacement for nightly recompute.
    \item \textbf{Q: What does \texttt{arg\_max(ingestion\_time(), *) by EventId} deduplicate?}\\
    A: At-least-once redeliveries: all physical copies of one logical event collapse to the latest-arrived row, keyed by the producer-stamped \texttt{EventId}.
    \item \textbf{Q: Why do unbounded lookback windows degrade query and materialization performance?}\\
    A: They abandon incremental state: every query and every view update re-scans ever-growing history instead of the recent slice, so cost grows with data age rather than with workload.
    \item \textbf{Q: How do late arrivals affect a time-windowed stream join?}\\
    A: Events arriving after the window closes never match, so joins undercount; the window must reflect the measured lateness distribution, and the policy must state how unmatched events are handled.
    \item \textbf{Q: Where must the dedup key be generated and enforced for end-to-end exactly-once under replay?}\\
    A: Generated at the producer before the first send, propagated untouched through Eventstream, and enforced at the first queryable boundary---KQL materialized view or Delta \texttt{MERGE}---so every downstream consumer reads clean data.
    \item \textbf{Q: What does consumer checkpointing protect against?}\\
    A: Restart-induced replay: recording the last processed offset or timestamp lets a resumed consumer continue rather than reprocess, which keeps both cost and duplicates bounded.
\end{enumerate}

\section{Further Reading}
Review the KQL documentation \cite{microsoftKQL}, materialized views \cite{microsoftKQLMV}, and the Eventhouse overview \cite{microsoftFabricEventhouse}. The ingestion path feeding these tables is Chapter~13's Eventstreams \cite{microsoftFabricEventstreams}; the reaction layer above them is Chapter~15's Data Activator \cite{microsoftDataActivator}.

\section*{Key Takeaways}
\begin{itemize}
    \item Eventhouses host KQL databases: append-first, time-native, policy-driven stores for telemetry.
    \item KQL composes through pipes: \texttt{where}, \texttt{project}, \texttt{summarize}, \texttt{bin()}---readable and time-obsessed.
    \item Materialized views are incremental state; dashboards should read views, not firehoses.
    \item \texttt{arg\_max(ingestion\_time(), *) by EventId} is the dedup contract made executable.
    \item \texttt{make-series} plus the \texttt{series\_*} family turns event streams into analyzable time series, including anomaly detection.
    \item Stream joins need time-windowed keys; late arrivals need a published policy, not silent revision.
    \item Unbounded lookback abandons incremental design---bound every query by time, always.
\end{itemize}

\section{Exercises}
\begin{enumerate}
    \item \textbf{(Conceptual)} Explain append-first design and why dedup happens at query time rather than ingestion time in a KQL database.
    \item \textbf{(Conceptual)} Contrast event time, ingestion time, and processing time; give one question each answers correctly.
    \item \textbf{(Hands-on)} Rewrite the lab's anomaly detection with \texttt{series\_decompose} and compare baseline/residual behavior to \texttt{series\_decompose\_anomalies}.
    \item \textbf{(Hands-on)} Build the materialized view version of the dedup pattern and measure the latency difference against query-time dedup at one day of volume.
    \item \textbf{(Design)} Write the lateness policy for the mobile-fleet dashboard: window sizes, provisional marking, and revision communication.
    \item \textbf{(Design)} Design the stream-to-dimension ``as-of'' enrichment when device metadata itself changes over time.
    \item \textbf{(Interview)} In five sentences, explain how materialized views, \texttt{arg\_max} dedup, and bounded lookback together deliver exactly-once \emph{effects} over at-least-once delivery.
    \item \textbf{(Operational)} Write the runbook for ``materialized view lagging behind ingestion''---detection, cause analysis, catch-up verification.
\end{enumerate}

\section{Capstone Project: Windowed Analytics over a Replay-Safe Stream}\label{sec:ch14-capstone}
\index{capstone project}\index{KQL!capstone}\index{materialized views!capstone}

\begin{capstonebox}
\textbf{Mission.} Deliver the analytical layer over Chapter~13's replay-safe stream: deduplicated base views, materialized per-device aggregates, a rolling-average anomaly detector, a time-windowed alert-correlation join, and a written lateness policy---with evidence that late and duplicated data cannot corrupt the published numbers.

\textbf{Estimated time.} 5--7 hours on top of the Chapter~13 estate.

\textbf{What you'll practice.}
\begin{itemize}
    \item Materialized view design for incremental aggregation.
    \item Anomaly detection with the \texttt{series\_*} family.
    \item Time-windowed stream joins and stream-to-dimension enrichment.
    \item Lateness policy as a published, tested contract.
\end{itemize}
\end{capstonebox}

\subsection*{Scenario}
Contoso Devices' support organization needs a live fleet-health view: per-site rolling temperatures, anomaly flags, and correlation between telemetry spikes and device alerts. The data platform team's requirement: yesterday's numbers must be stable by 09:00, with a documented policy covering the fleet's nine-hour offline backhauls.

\subsection*{Requirements}
\begin{itemize}
    \item Deduped base consumption (view or materialized view) enforcing \texttt{arg\_max} on \texttt{EventId}.
    \item Materialized view of per-device, per-site five-minute means.
    \item The rolling-average anomaly query parameterized by threshold.
    \item The alert-correlation join with a justified time window.
    \item The lateness policy document: lookback, provisional marking, stabilization time.
    \item Evidence runs: injected duplicates and a simulated nine-hour backhaul, both absorbed correctly.
\end{itemize}

\subsection*{Implementation Steps}
\begin{enumerate}
    \item Build the deduped base and the materialized aggregates.
    \item Implement and tune the anomaly query against injected spikes.
    \item Implement the correlation join; justify the 300-second window or revise it with data.
    \item Write the lateness policy; implement provisional marking in the consumption queries.
    \item Run both evidence scenarios and capture query outputs.
\end{enumerate}

\subsection*{Deliverables}
\begin{itemize}
    \item All KQL artifacts (views, materialized views, functions) as versioned scripts.
    \item Anomaly and correlation query results on the evidence data.
    \item The lateness policy document.
    \item Evidence outputs proving duplicate and backhaul absorption.
\end{itemize}

\subsection*{Acceptance Criteria}
\begin{itemize}
    \item[\(\square\)] Injected duplicates never appear in any published aggregate.
    \item[\(\square\)] The simulated backhaul revises only provisional windows, never stabilized ones.
    \item[\(\square\)] Every consumption query is time-bounded; none scans full history.
    \item[\(\square\)] The anomaly detector fires on the synthetic spike and stays quiet on noise.
    \item[\(\square\)] The correlation join's window is justified with measured latency data.
    \item[\(\square\)] The lateness policy is written, reviewed, and referenced by the dashboard contract.
\end{itemize}

\subsection*{Stretch Goals}
\begin{itemize}
    \item Add a forecast (\texttt{series\_decompose\_forecast}) for next-hour site temperature and track its error.
    \item Implement as-of dimension joins for device firmware version at event time.
    \item Compare materialized-view cost versus on-demand aggregation at seven days of volume; publish the break-even.
    \item Feed the anomaly flags to a Reflex item (previews Chapter~15) and measure detection-to-alert latency.
\end{itemize}
